% Part: computability
% Chapter: computability-theory
% Section: complete-ce-sets

\documentclass[../../../include/open-logic-section]{subfiles}

\begin{document}

\olfileid{cmp}{thy}{cce}
\olsection{પૂર્ણ સંગણકીય રીતે પરિગણનીય ગણો}

\begin{defn}
ગણ $A$ને (બહુ-એક ન્યૂનીકરણક્ષમતાની દૃષ્ટિએ)
\emph{પૂર્ણ !!{computably enumerable} ગણ} ત્યારે કહીએ જ્યારે
\begin{enumerate}
\item $A$ સંગણકીય રીતે પરિગણનીય હોય અને
\item બીજા દરેક સંગણકીય રીતે પરિગણનીય ગણ $B$ માટે $B \leq_m A$ હોય.
\end{enumerate}
\end{defn}

બીજા શબ્દોમાં, પૂર્ણ સંગણકીય રીતે પરિગણનીય ગણો શક્ય
સૌથી ``કઠિન'' એવા પરિગણનીય ગણો છે. તેઓ \emph{કોઈ પણ}
સંગણકીય રીતે પરિગણનીય ગણ વિશેના પ્રશ્નોના જવાબ મેળવવા દે છે.

\begin{thm}
$K$, $K_0$ અને આગળ વ્યાખ્યાયિત થનારો $K_1$ બધા પૂર્ણ
સંગણકીય રીતે પરિગણનીય ગણો છે.
\end{thm}

\begin{proof}
$K_0$ પૂર્ણ છે તે જોવા કોઈ પણ સંગણકીય રીતે પરિગણનીય ગણ
$B$ લો. ત્યારે કોઈ સૂચક $e$ માટે
\[
B = W_e = \Setabs{x}{\cfind{e}(x) \fdefined}.
\]
$f$ એ $f(x) = \tuple{e, x}$ વિધેય હોય. ત્યારે દરેક પ્રાકૃતિક
સંખ્યા $x$ માટે $x \in B$ ત્યારે અને માત્ર ત્યારે જ્યારે
$f(x) \in K_0$. બીજા શબ્દોમાં, $f$, $B$નું~$K_0$માં
ન્યૂનીકરણ કરે છે.

$K_1$ પૂર્ણ છે તે જોવા નોંધો કે \olref[k1]{prop:k1}ની
સાબિતીમાં આપણે $K_0$નું તેમાં ન્યૂનીકરણ કર્યું છે. તેથી
\olref[ppr]{prop:trans-red} મુજબ કોઈ પણ સંગણકીય રીતે
પરિગણનીય ગણને~$K_1$માં ઘટાડી શકાય છે.

$K_0$નું $K$માં પણ લગભગ આ જ રીતે ન્યૂનીકરણ થાય છે:
ઉપરના કાર્યક્રમ-સૂચકના બાંધકામની જેમ, ઇનપુટને અવગણીને
આપેલી સંગણના ચલાવતો સૂચક બનાવો.
\sourcecorrection{OLCT-009}{સ્થિર અંગ્રેજી મૂળ છેલ્લે
$K$નું $K_0$માં ન્યૂનીકરણ કહે છે. તે વિરુદ્ધ દિશા છે અને
$K$ પૂર્ણ છે એ દાવાને સાબિત કરતી નથી. જરૂરી દિશા
$K_0$માંથી $K$ તરફ છે.}
\end{proof}

\begin{prob}
$K$નું $K_0$માં ન્યૂનીકરણ આપો.
\end{prob}

\begin{digress}
આમ, અત્યાર સુધી આપણે જોયેલા સંગણકીય રીતે પરિગણનીય ગણોના
બધાં ઉદાહરણો કાં તો સંગણનીય છે અથવા પૂર્ણ છે. આ વિચિત્ર
લાગવું જોઈએ!{} શું એવા સંગણકીય રીતે પરિગણનીય ગણો છે જે
સંગણનીય પણ નથી અને પૂર્ણ પણ નથી? જવાબ હા છે, પરંતુ
ફ્રિડબર્ગ અને મુચનિકે સ્વતંત્ર રીતે ૧૯૫૦ના દાયકાના મધ્યમાં
આ સ્થાપ્યું ત્યાં સુધી આ જાણીતું નહોતું.
\end{digress}

\end{document}
